\input{preamble.tex}

\title{Отчет по лабораторной работе № 4 \\ Первоначальное конфигурирование сети}
\author{Сущенко Алина \\ НПИбд-01-23}
\date{2026}

\begin{document}

\maketitle
\newpage

\tableofcontents
\newpage

\section{Цель работы}
Провести подготовительную работу по первоначальной настройке коммутаторов сети.

\section{Задание}
Сделать первоначальную настройку коммутаторов сети, представленной на схеме L1. Под первоначальной настройкой понимается указание имени устройства, его IP-адреса, настройка доступа по паролю к виртуальным терминалам и консоли, настройка удалённого доступа к устройству по ssh. При выполнении работы необходимо учитывать соглашение об именовании.

\section{Выполнение работы}

\subsection{Размещение коммутаторов и оконечных устройств}
В логической рабочей области \texttt{Packet Tracer} разместили коммутаторы и оконечные устройства согласно схеме сети L1. Выдали им названия в соответствии с соглашением об именовании \texttt{ansusenko} (Рис. \ref{01}):

\begin{center}
    \centering
    \includegraphics[width=0.9\textwidth]{../images/01.png}
    \captionof{figure}{Размещение коммутаторов и оконечных устройств согласно схеме сети L1}
    \label{01}
\end{center}

\subsection{Настройка коммутатора msk-ansusenko-sw-01}
Произвели первоначальную настройку первого коммутатора. Вошли в режим глобальной конфигурации, задали имя устройства, настроили VLAN 2 с IP-адресом, указали шлюз по умолчанию (Рис. \ref{02}):

\begin{center}
    \centering
    \includegraphics[width=0.9\textwidth]{../images/02.png}
    \captionof{figure}{Настройка имени, IP-адреса и шлюза на msk-ansusenko-sw-01}
    \label{02}
\end{center}

Настроили пароли для доступа к виртуальным терминалам (VTY) и консоли, установили секретный пароль для привилегированного режима, включили шифрование паролей
\begin{itemize}
    \item Пароль на линии VTY 0-4: \texttt{cisco}
    \item Пароль на консоли: \texttt{cisco}
    \item Секретный пароль enable: \texttt{cisco}
    \item Включено шифрование паролей командой \texttt{service password-encryption}
\end{itemize}

Проверим как настроена конфигурация на \texttt{msk-ansusenko-sw-01} (Рис. \ref{03}-\ref{05}):

\begin{center}
    \centering
    \includegraphics[width=0.9\textwidth]{../images/03.png}
    \captionof{figure}{Просмотр текущей конфигурации msk-ansusenko-sw-01}
    \label{03}
\end{center}

\begin{center}
    \centering
    \includegraphics[width=0.9\textwidth]{../images/04.png}
    \captionof{figure}{Продолжение конфигурации msk-ansusenko-sw-01 1}
    \label{04}
\end{center}

\begin{center}
    \centering
    \includegraphics[width=0.9\textwidth]{../images/05.png}
    \captionof{figure}{Продолжение конфигурации msk-ansusenko-sw-01 2}
    \label{05}
\end{center}

\subsection{Настройка коммутатора msk-ansusenko-sw-1}
Произвели настройку второго коммутатора с именем \texttt{msk-ansusenko-sw-1}, присвоили ему IP-адрес 10.128.1.6 (Рис. \ref{06}):

\begin{center}
    \centering
    \includegraphics[width=0.9\textwidth]{../images/06.png}
    \captionof{figure}{Настройка коммутатора msk-ansusenko-sw-1}
    \label{06}
\end{center}

После настройки проверили текущую конфигурацию и состояние интерфейсов (Рис. \ref{07}):

\begin{center}
    \centering
    \includegraphics[width=0.9\textwidth]{../images/07.png}
    \captionof{figure}{Просмотр running-config и состояния интерфейсов на msk-ansusenko-sw-1}
    \label{07}
\end{center}

Просмотрели информацию о VLAN и транках (Рис. \ref{08}):

\begin{center}
    \centering
    \includegraphics[width=0.9\textwidth]{../images/08.png}
    \captionof{figure}{Информация о VLAN и транках на msk-ansusenko-sw-1}
    \label{08}
\end{center}

\subsection{Настройка коммутатора msk-ansusenko-sw-2}
Настроили коммутатор \texttt{msk-ansusenko-sw-2} с IP-адресом 10.128.1.3 (Рис. \ref{09}):

\begin{center}
    \centering
    \includegraphics[width=0.9\textwidth]{../images/09.png}
    \captionof{figure}{Настройка коммутатора msk-ansusenko-sw-2}
    \label{09}
\end{center}

Проверили конфигурацию и состояние интерфейсов (Рис. \ref{10}):

\begin{center}
    \centering
    \includegraphics[width=0.9\textwidth]{../images/10.png}
    \captionof{figure}{Просмотр конфигурации и интерфейсов на msk-ansusenko-sw-2}
    \label{10}
\end{center}

Просмотрели информацию о VLAN и транках на данном коммутаторе (Рис. \ref{11}):

\begin{center}
    \centering
    \includegraphics[width=0.9\textwidth]{../images/11.png}
    \captionof{figure}{Информация о VLAN и транках на msk-ansusenko-sw-2}
    \label{11}
\end{center}

\subsection{Настройка коммутатора msk-ansusenko-sw-3}
Настроили коммутатор \texttt{msk-ansusenko-sw-3} и проверили информацию о VLAN (Рис. \ref{12}-\ref{14}):

\begin{center}
    \centering
    \includegraphics[width=0.9\textwidth]{../images/12.png}
    \captionof{figure}{Настройка msk-ansusenko-sw-3}
    \label{12}
\end{center}

\begin{center}
    \centering
    \includegraphics[width=0.9\textwidth]{../images/13.png}
    \captionof{figure}{Информация об msk-ansusenko-sw-3}
    \label{13}
\end{center}

\begin{center}
    \centering
    \includegraphics[width=0.9\textwidth]{../images/14.png}
    \captionof{figure}{Информация о VLAN на msk-ansusenko-sw-3}
    \label{14}
\end{center}

\subsection{Настройка коммутатора msk-ansusenko-sw-4}
Настроили коммутатор \texttt{msk-ansusenko-sw-4} (Рис. \ref{15}):

\begin{center}
    \centering
    \includegraphics[width=0.9\textwidth]{../images/15.png}
    \captionof{figure}{Настройка коммутатора msk-ansusenko-sw-4}
    \label{15}
\end{center}

Просмотрели информацию о VLAN на последнем коммутаторе (Рис. \ref{16}-\ref{17}):

\begin{center}
    \centering
    \includegraphics[width=0.9\textwidth]{../images/16.png}
    \captionof{figure}{Информация на msk-ansusenko-sw-4}
    \label{16}
\end{center}

\begin{center}
    \centering
    \includegraphics[width=0.9\textwidth]{../images/17.png}
    \captionof{figure}{Информация о VLAN на msk-ansusenko-sw-4}
    \label{17}
\end{center}

\subsection{Сводная информация по настройкам}
В процессе выполнения работы все коммутаторы были настроены в соответствии с планом IP-адресации (Таблица \ref{tab:ip}):

\begin{table}[h]
    \centering
    \caption{Назначенные IP-адреса коммутаторов}
    \begin{tabular}{|c|c|c|}
        \hline
        \textbf{Устройство} & \textbf{Имя} & \textbf{IP-адрес} \\
        \hline
        Коммутатор 1 & msk-ansusenko-sw-01 & 10.128.1.2 \\
        \hline
        Коммутатор 2 & msk-ansusenko-sw-1 & 10.128.1.6 \\
        \hline
        Коммутатор 3 & msk-ansusenko-sw-2 & 10.128.1.3 \\
        \hline
        Коммутатор 4 & msk-ansusenko-sw-3 & 10.128.1.4 \\
        \hline
        Коммутатор 4 & msk-ansusenko-sw-4 & 10.128.1.5 \\
    \end{tabular}
    \label{tab:ip}
\end{table}

\section{Контрольные вопросы}

\subsection{Вопрос 1: При помощи каких команд можно посмотреть конфигурацию сетевого оборудования?}
Для просмотра конфигурации сетевого оборудования Cisco используются следующие команды:

\begin{itemize}
    \item \texttt{show running-config} (или сокращенно \texttt{show run}) - отображает текущую конфигурацию, находящуюся в оперативной памяти (RAM) и используемую устройством в данный момент.
    
    \item \texttt{show startup-config} (или \texttt{show start}) - отображает стартовую конфигурацию, сохраненную в энергонезависимой памяти (NVRAM), которая загружается при включении устройства.
    
    \item \texttt{show ip interface brief} - показывает краткую информацию по всем интерфейсам, включая IP-адреса и статус (up/down).
    
    \item \texttt{show vlan} - отображает информацию о настроенных VLAN и принадлежности портов к ним (как показано на Рис. \ref{fig:08}, \ref{fig:11}, \ref{fig:14}, \ref{fig:17}).
    
    \item \texttt{show interfaces trunk} - показывает информацию о транковых портах и разрешенных VLAN на них (Рис. \ref{fig:08}, \ref{fig:11}).
    
    \item \texttt{show ssh} - отображает информацию о настроенном SSH и активных соединениях.
\end{itemize}

\subsection{Вопрос 2: При помощи каких команд можно посмотреть стартовый конфигурационный файл оборудования?}
Стартовый конфигурационный файл (startup-config) просматривается командой:
\begin{verbatim}
    show startup-config
\end{verbatim}
или сокращенно:
\begin{verbatim}
    show start
\end{verbatim}
Эта конфигурация хранится в NVRAM и загружается при включении или перезагрузке устройства. В отличие от running-config, startup-config сохраняется после выключения питания.

\subsection{Вопрос 3: При помощи каких команд можно экспортировать конфигурационный файл оборудования?}
Для экспорта (сохранения) конфигурации используются следующие команды:

\begin{enumerate}
    \item \textbf{Сохранение в NVRAM устройства:}
    \begin{verbatim}
        write memory
    \end{verbatim}
    или
    \begin{verbatim}
        copy running-config startup-config
    \end{verbatim}
    Эта команда сохраняет текущую конфигурацию в энергонезависимую память, чтобы она загружалась при перезагрузке (использовалась в работе, см. Рис. \ref{fig:06}, \ref{fig:15}).
    
    \item \textbf{Экспорт на внешний TFTP-сервер:}
    \begin{verbatim}
        copy running-config tftp:
    \end{verbatim}
    После ввода команды потребуется указать IP-адрес TFTP-сервера и имя файла для сохранения.
    
    \item \textbf{Экспорт на FTP/SCP сервер:}
    \begin{verbatim}
        copy running-config ftp:
        copy running-config scp:
    \end{verbatim}
\end{enumerate}

\subsection{Вопрос 4: При помощи каких команд можно импортировать конфигурационный файл оборудования?}
Для импорта (загрузки) конфигурации используются следующие команды:

\begin{enumerate}
    \item \textbf{Загрузка из NVRAM в оперативную память:}
    \begin{verbatim}
        copy startup-config running-config
    \end{verbatim}
    Эта команда загружает сохраненную конфигурацию в текущую (используется для отката изменений).
    
    \item \textbf{Загрузка с внешнего TFTP-сервера:}
    \begin{verbatim}
        copy tftp: running-config
        copy tftp: startup-config
    \end{verbatim}
    В первом случае конфигурация загружается в текущую сессию, во втором — сразу сохраняется как стартовая.
    
    \item \textbf{Загрузка с FTP/SCP сервера:}
    \begin{verbatim}
        copy ftp: running-config
        copy scp: startup-config
    \end{verbatim}
    
    \item \textbf{Перезагрузка устройства для применения стартовой конфигурации:}
    \begin{verbatim}
        reload
    \end{verbatim}
    При перезагрузке устройство загружает конфигурацию из startup-config.
\end{enumerate}

\section{Выводы}
В результате выполнения лабораторной работы приобрели практические навыки по первоначальной настройке коммутаторов Cisco в среде Packet Tracer.

\end{document}
